\subsection{Building Blocks of Variational Algorithms}\label{sec:building-blocks-of-variational-algorithms}

\subsubsection{Cost Function}\label{sec:cost-function}

A \definition{Cost Function} is a \textbf{numerical measure of the quality of a candidate solution}. It maps the quantum state generated by the ansatz to a classical number that the optimizer can evaluate. The cost function is designed so that its optimum value corresponds to the solution of the original problem.

\highspace
\begin{flushleft}
    \textcolor{Green3}{\faIcon{book} \textbf{The Cost Function Defines the Problem}}
\end{flushleft}
In traditional algorithms, it is the circuit that defines the problem. In contrast, in variational quantum algorithms, it is the cost function that defines the problem. The quantum circuit (ansatz) is just a tool to explore the solution space, while the \textbf{cost function encodes the specific problem we want to solve}. For example, in VQE, the cost function is the expectation value of the Hamiltonian, which we want to minimize to find the ground state energy.

\highspace
\begin{flushleft}
    \textcolor{Green3}{\faIcon{check-circle} \textbf{What Makes a Good Cost Function?}}
\end{flushleft}
The \textbf{minimum of the cost function must correspond to the solution of the problem}. If this is not true, the optimizer may find a minimum that has nothing to do with the problem we want to solve. Therefore, the cost function must be carefully designed to ensure that its minimum corresponds to the desired solution. Additionally, the cost function should be efficiently computable on a quantum computer, as it will be evaluated many times during the optimization process.

\highspace
\begin{flushleft}
    \textcolor{Green3}{\faIcon{exchange-alt} \textbf{The Cost Function Connects Quantum and Classical Worlds}}
\end{flushleft}
The cost function serves as the \textbf{bridge between the quantum and classical parts of the algorithm}. The quantum computer produces measurements; the cost function takes these measurements and produces a classical number that the optimizer can use. This is a crucial point: the \hl{optimizer does not need to understand quantum mechanics; it only needs to know how to evaluate the cost function based on the measurement outcomes}. The \textbf{cost function translates quantum information into a form that classical optimization algorithms can work with}.

\highspace
In summary, the \definition{Cost Function} is the quantity optimized by a Variational Quantum Algorithm (VQA). It evaluates the quality of the quantum state generated by the ansatz and provides a numerical score to the classical optimizer. The cost function must be designed so that its minimum (or maximum, depending on the formulation) corresponds to the solution of the problem. Therefore, the cost function effectively defines the problem being solved by the variational algorithm. For this reason, the cost function is the \textbf{most important part}. \hl{Even if we have a fantastic optimizer and a powerful ansatz, if the cost function is poorly designed, we will solve the wrong problem or fail to find the correct solution.}